% chapters/09/section-04.tex
\section{Menguji Setiap Keadaan dan Perpindahannya}
\label{sec:menguji-keadaan}

Mesin keadaan sulit diperiksa dengan cara biasa, sebab perilakunya bergantung
pada posisi pemain dan pada waktu. Bagian ini membahas cara mengujinya, sesuai
indikator \ind{c} pada bab ini.

\subsection{Mengapa Pengujiannya Sulit}

Bayangkan Anda ingin memeriksa keadaan menyerang. Anda harus mendekatkan pemain
sampai sedekat mungkin, menunggu NPC menyadarinya, lalu menunggu NPC mendekat.
Setiap percobaan menuntut puluhan detik, dan hasilnya bergantung pada
keterampilan Anda mengendalikan pemain. Cara itu jelas tidak memadai untuk
memeriksa tujuh perpindahan yang berbeda.

Tiga cara berikut dipakai bersama-sama pada bab ini.

\subsection{Cara Pertama: Mencatat Perpindahan beserta Bingkainya}

Program mencatat setiap perpindahan keadaan beserta nomor bingkai ketika
perpindahan itu terjadi. Catatan itu ditampilkan pada jendela selama permainan
berlangsung dan dicetak pada terminal ketika program berakhir.

\begin{kodekeluaran}
    0  berdiam
   91  berpatroli
  151  mengejar
  217  menyerang
  421  mengejar
  482  berdiam
  571  berpatroli
\end{kodekeluaran}

Catatan itu jauh lebih berguna daripada tangkapan layar. Ia memperlihatkan
bukan hanya keadaan yang sedang aktif, melainkan urutan dan lamanya setiap
keadaan. Dari catatan di atas, misalnya, tampak bahwa NPC berdiam selama 91
bingkai, yaitu sekitar 1,5 detik, sesuai nilai pewaktu yang ditetapkan.

\subsection{Cara Kedua: Skenario yang Dapat Diulang}

Catatan perpindahan menjadi alat pengujian yang sesungguhnya apabila
permainannya dapat dijalankan dengan skenario tetap. Program bab ini
menyediakan pilihan \code{--scenario} yang menggerakkan pemain menurut jalur
yang telah ditentukan, sehingga seluruh perpindahan terjadi dalam urutan yang
dapat diperkirakan.

Skenario itu menempatkan pemain pada tiga posisi, masing-masing selama jangka
waktu yang cukup panjang, yaitu di luar jangkauan persepsi, di dalam jangkauan
persepsi, lalu di luar jangkauan persepsi kembali. Dengan begitu, ketujuh
perpindahan pada Tabel~\ref{tab:tujuh-perpindahan} terjadi seluruhnya.

\begin{catatabox}
Skenario itu menuntut pemain berdiri pada jarak yang
\textbf{cukup jauh} ketika menghilang. Apabila jaraknya kurang jauh, NPC akan
menemukan pemain kembali sebelum pewaktu tenggangnya habis, sebab NPC tetap
bergerak mengejar selama pewaktu itu berjalan. Jarak yang diperlukan kira-kira
sama dengan jari-jari persepsi ditambah kecepatan NPC dikalikan lama pewaktu
tenggang. Pada nilai yang dipakai bab ini, jaraknya sekitar 410 piksel.
\end{catatabox}

\subsection{Cara Ketiga: Memeriksa Keabsahan Perpindahan}

Cara kedua menuntut urutan yang tepat, sehingga hasilnya peka terhadap
perbedaan kecepatan komputer. Karena itu, program bab ini juga memeriksa hal
yang lebih mendasar: apakah setiap perpindahan yang terjadi termasuk
perpindahan yang sah.

\begin{kodekeluaran}
LEGAL_TRANSITIONS = {
    IDLE_STATE:   {PATROL_STATE, CHASE_STATE},
    PATROL_STATE: {IDLE_STATE, CHASE_STATE},
    CHASE_STATE:  {ATTACK_STATE, IDLE_STATE},
    ATTACK_STATE: {CHASE_STATE},
}
\end{kodekeluaran}

Tabel itu adalah diagram pada Gambar~\ref{fig:diagram-keadaan} yang ditulis
kembali dalam bentuk yang dapat diperiksa program. Setiap pasangan keadaan yang
berurutan pada catatan perpindahan diperiksa terhadap tabel itu.

Pemeriksaan semacam ini disebut pemeriksaan \textbf{sifat yang harus selalu
berlaku}. Ia tidak menuntut urutan tertentu, sehingga tahan terhadap perbedaan
waktu antar-komputer, tetapi tetap mampu menangkap perpindahan yang keliru.
Pada mesin keadaan yang lebih rumit, cara inilah yang paling berguna.

\subsection{Memandang Keadaan sebagai Garis Waktu}

Catatan perpindahan dapat pula digambarkan sebagai garis waktu, sehingga lama
setiap keadaan terlihat sekaligus. Gambar~\ref{fig:garis-waktu-keadaan}
memperlihatkannya untuk skenario pengujian di atas.

\begin{figure}[htbp]
\centering
% chapters/09/figures/garis-waktu-keadaan.tikz
% Garis waktu keadaan NPC pada skenario pengujian, memakai catatan perpindahan
% yang sesungguhnya: setiap pita menyatakan lama sebuah keadaan berlangsung.
\begin{tikzpicture}[
  pita/.style={rectangle, draw=ubLine, thick, minimum height=9mm,
               align=center, font=\scriptsize},
  label/.style={font=\scriptsize, align=center},
  penanda/.style={draw=ubLine, thin}
]

% Pita keadaan. Lebar sebanding dengan lama keadaan berlangsung.
\node[pita, fill=ubGray, minimum width=1.42cm] (a) at (0.71,0) {berdiam};
\node[pita, fill=ubGreenBg, minimum width=0.93cm, right=0cm of a] (b) {patroli};
\node[pita, fill=ubRedBg, minimum width=1.03cm, right=0cm of b] (c) {kejar};
\node[pita, fill=ubWarnBg, minimum width=3.17cm, right=0cm of c] (d) {menyerang};
\node[pita, fill=ubRedBg, minimum width=0.95cm, right=0cm of d] (e) {kejar};
\node[pita, fill=ubGray, minimum width=1.42cm, right=0cm of e] (f) {berdiam};
\node[pita, fill=ubGreenBg, minimum width=5.09cm, right=0cm of f] (g) {berpatroli};

% Nomor bingkai pada setiap batas keadaan.
\foreach \titik/\bingkai in {a.west/0, a.east/91, b.east/151, c.east/217,
                             d.east/421, e.east/482, f.east/573, g.east/900} {
  \draw[penanda] (\titik) ++(0,0.75) -- ++(0,0.25);
  \node[label, anchor=south] at ($(\titik) + (0,1.0)$) {\bingkai};
}

% Sumbu waktu.
\draw[-{Stealth[length=2mm]}, thick, draw=ubLine] (-0.3,-0.9) -- (14.6,-0.9);
\node[label, anchor=north west] at (0,-1.1)
  {bingkai (900 bingkai setara 15 detik pada 60 bingkai per detik)};

\node[label, anchor=north west, text width=140mm] at (-0.3,-1.9)
  {Pita berwarna menyatakan keadaan yang sedang aktif, dan lebarnya sebanding
   dengan lama keadaan itu berlangsung. Perhatikan bahwa keadaan menyerang
   berlangsung paling lama, yaitu 204 bingkai, sebab pemain berdiri di dalam
   jangkauan serang selama 3,4 detik.};
\end{tikzpicture}

\caption{Garis waktu keadaan NPC pada skenario pengujian}
\label{fig:garis-waktu-keadaan}
\end{figure}

Gambar itu memperlihatkan hal yang tidak terlihat pada catatan angka, misalnya
bahwa keadaan menyerang berlangsung paling lama, yaitu 204 bingkai. Angka itu
dapat diperiksa dengan perhitungan: pemain berada di dalam jangkauan serang
selama 3,4 detik, sehingga jumlah serangan seharusnya adalah 3,4 dibagi 0,8
ditambah satu, yaitu sekitar lima serangan. Program mencatat lima serangan.

\begin{catatabox}
Apabila hasil pengujian tidak sesuai dengan dugaan Anda, periksalah lebih
dahulu apakah dugaan Anda yang keliru, bukan programnya. Pengalaman penulis
selama menyusun bab ini persis demikian: skenario pengujian semula dirancang
dengan jarak yang tidak cukup jauh, sehingga NPC menemukan pemain kembali dan
urutan keadaannya berbeda dari dugaan. Programnya benar; rancangan
skenarionya yang perlu diperbaiki.
\end{catatabox}
